
\documentclass{article}
\usepackage{colortbl}
\usepackage{makecell}
\usepackage{multirow}
\usepackage{supertabular}

\begin{document}

\newcounter{utterance}

\centering \large Interaction Transcript for game `dond', experiment `semi', episode 2 with nemotron{-}3{-}ultra.
\vspace{24pt}

{ \footnotesize  \setcounter{utterance}{1}
\setlength{\tabcolsep}{0pt}
\begin{supertabular}{c@{$\;$}|p{.15\linewidth}@{}p{.15\linewidth}p{.15\linewidth}p{.15\linewidth}p{.15\linewidth}p{.15\linewidth}}
    \# & $\;$A & \multicolumn{4}{c}{Game Master} & $\;\:$B\\
    \hline

    \theutterance \stepcounter{utterance}  
    & & \multicolumn{4}{p{0.6\linewidth}}{
        \cellcolor[rgb]{0.9,0.9,0.9}{
            \makecell[{{p{\linewidth}}}]{
                \texttt{\tiny{[P1$\langle$GM]}}
                \texttt{Παίζετε ένα συνεργατικό παιχνίδι διαπραγμάτευσης στο οποίο πρέπει να συμφωνήσετε με έναν άλλο παίκτη σχετικά με τον τρόπο διαίρεσης ενός συνόλου αντικειμένων μεταξύ σας.} \\
\\ 
\texttt{Κανόνες:} \\
\texttt{(a) Εσείς και ο άλλος παίκτης λαμβάνετε ένα σύνολο αντικειμένων. Σας δίνεται επίσης μια μυστική συνάρτηση τιμής που αντιπροσωπεύει πόσο εκτιμάτε κάθε τύπο αντικειμένου.} \\
\texttt{(b) Ανταλλάσσετε μηνύματα με τον άλλο παίκτη για να συμφωνήσετε σχετικά με το ποιος παίρνει ποια αντικείμενα. Μπορείτε να στείλετε το πολύ 5 μηνύματα ο καθένας ή να τερματίσετε νωρίς κάνοντας μια μυστική πρόταση ανά πάσα στιγμή.} \\
\texttt{(c) Σας ζητείται να υποβάλετε μια μυστική πρόταση που υποδεικνύει τα αντικείμενα που θέλετε, μορφοποιημένη σε τετράγωνες αγκύλες ως εξής: "{[}Πρόταση: <αριθμός> <όνομα αντικειμένου>, <αριθμός> <όνομα αντικειμένου>, <...>{]}"} \\
\texttt{(d) Εάν οι προτάσεις σας είναι συμπληρωματικές, δηλαδή υπάρχουν αρκετά αντικείμενα για να ικανοποιήσουν και τις δύο προτάσεις, κάθε παίκτης λαμβάνει βαθμολογία με βάση το άθροισμα των τιμών των αντικειμένων που έλαβε. Διαφορετικά, και οι δύο παίρνετε μηδέν πόντους.} \\
\texttt{(e) Ο στόχος σας είναι να μεγιστοποιήσετε το άθροισμα της βαθμολογίας σας και της βαθμολογίας του άλλου παίκτη.} \\
\\ 
\texttt{Ας ξεκινήσουμε.} \\
\\ 
\texttt{Το σύνολο των διαθέσιμων αντικειμένων είναι:} \\
\\ 
\texttt{1 τσάντα, 3 λάμπες, 1 λουλούδι} \\
\\ 
\texttt{Η μυστική συνάρτηση τιμής σας είναι:} \\
\\ 
\texttt{τσάντα: 3, λάμπα: 0, λουλούδι: 7} \\
\\ 
\texttt{ΣΗΜΑΝΤΙΚΟ: Τα μηνύματά σας, εκτός εάν είναι μυστική πρόταση, μεταδίδονται απευθείας στον άλλο παίκτη, επομένως μην συμπεριλάβετε καμία απάντηση στους κανόνες ή κείμενο που ανακοινώνει το μήνυμά σας. Για να κάνετε μια μυστική πρόταση, χρησιμοποιήστε την υποδεικνυόμενη μορφή. Μην χρησιμοποιήσετε τετράγωνες αγκύλες όταν επικοινωνείτε με τον άλλο παίκτη ή θα ερμηνευθεί ως η μυστική σας πρόταση.} \\
            }
        }
    }
    & \\ \\

    \theutterance \stepcounter{utterance}  
    & \multicolumn{4}{p{0.6\linewidth}}{
        \cellcolor[rgb]{0.9,0.9,0.9}{
            \makecell[{{p{\linewidth}}}]{
                \texttt{\tiny{[P1$\rangle$GM]}}
                \texttt{Γεια σου! Έχουμε έναν πολύ απλό διαχωρισμό εδώ: εγώ αξιολογώ μόνο το λουλούδι (7 πόντους) και την τσάντα (3 πόντους), ενώ οι λάμπες δεν μου δίνουν τίποτα.} \\
\\ 
\texttt{Ποια είναι η τιμή σου για κάθε αντικείμενο; Αν τις μοιραστείς, μπορούμε να βρούμε την βέλτιστη συμφωνία γρήγορα.} \\
            }
        }
    }
    & & \\ \\

    \theutterance \stepcounter{utterance}  
    & & & \multicolumn{2}{p{0.3\linewidth}}{
        \cellcolor[rgb]{0.9,0.9,0.9}{
            \makecell[{{p{\linewidth}}}]{
                \texttt{\tiny{[GM$|$GM]}}
                \texttt{continue} \\
            }
        }
    }
    & & \\ \\

    \theutterance \stepcounter{utterance}  
    & & & \multicolumn{2}{p{0.3\linewidth}}{
        \cellcolor[rgb]{0.9,0.9,0.9}{
            \makecell[{{p{\linewidth}}}]{
                \texttt{\tiny{[GM$|$GM]}}
                \texttt{Γεια σου! Έχουμε έναν πολύ απλό διαχωρισμό εδώ: εγώ αξιολογώ μόνο το λουλούδι (7 πόντους) και την τσάντα (3 πόντους), ενώ οι λάμπες δεν μου δίνουν τίποτα.} \\
\\ 
\texttt{Ποια είναι η τιμή σου για κάθε αντικείμενο; Αν τις μοιραστείς, μπορούμε να βρούμε την βέλτιστη συμφωνία γρήγορα.} \\
            }
        }
    }
    & & \\ \\

    \theutterance \stepcounter{utterance}  
    & & \multicolumn{4}{p{0.6\linewidth}}{
        \cellcolor[rgb]{0.9,0.9,0.9}{
            \makecell[{{p{\linewidth}}}]{
                \texttt{\tiny{[GM$\rangle$P2]}}
                \texttt{Παίζετε ένα συνεργατικό παιχνίδι διαπραγμάτευσης στο οποίο πρέπει να συμφωνήσετε με έναν άλλο παίκτη σχετικά με τον τρόπο διαίρεσης ενός συνόλου αντικειμένων μεταξύ σας.} \\
\\ 
\texttt{Κανόνες:} \\
\texttt{(a) Εσείς και ο άλλος παίκτης λαμβάνετε ένα σύνολο αντικειμένων. Σας δίνεται επίσης μια μυστική συνάρτηση τιμής που αντιπροσωπεύει πόσο εκτιμάτε κάθε τύπο αντικειμένου.} \\
\texttt{(b) Ανταλλάσσετε μηνύματα με τον άλλο παίκτη για να συμφωνήσετε σχετικά με το ποιος παίρνει ποια αντικείμενα. Μπορείτε να στείλετε το πολύ 5 μηνύματα ο καθένας ή να τερματίσετε νωρίς κάνοντας μια μυστική πρόταση ανά πάσα στιγμή.} \\
\texttt{(c) Σας ζητείται να υποβάλετε μια μυστική πρόταση που υποδεικνύει τα αντικείμενα που θέλετε, μορφοποιημένη σε τετράγωνες αγκύλες ως εξής: "{[}Πρόταση: <αριθμός> <όνομα αντικειμένου>, <αριθμός> <όνομα αντικειμένου>, <...>{]}"} \\
\texttt{(d) Εάν οι προτάσεις σας είναι συμπληρωματικές, δηλαδή υπάρχουν αρκετά αντικείμενα για να ικανοποιήσουν και τις δύο προτάσεις, κάθε παίκτης λαμβάνει βαθμολογία με βάση το άθροισμα των τιμών των αντικειμένων που έλαβε. Διαφορετικά, και οι δύο παίρνετε μηδέν πόντους.} \\
\texttt{(e) Ο στόχος σας είναι να μεγιστοποιήσετε το άθροισμα της βαθμολογίας σας και της βαθμολογίας του άλλου παίκτη.} \\
\\ 
\texttt{Ας ξεκινήσουμε.} \\
\\ 
\texttt{Το σύνολο των διαθέσιμων αντικειμένων είναι:} \\
\\ 
\texttt{1 τσάντα, 3 λάμπες, 1 λουλούδι} \\
\\ 
\texttt{Η μυστική συνάρτηση τιμής σας είναι:} \\
\\ 
\texttt{τσάντα: 1, λάμπα: 2, λουλούδι: 3} \\
\\ 
\texttt{ΣΗΜΑΝΤΙΚΟ: Τα μηνύματά σας, εκτός εάν είναι μυστική πρόταση, μεταδίδονται απευθείας στον άλλο παίκτη, επομένως μην συμπεριλάβετε καμία απάντηση στους κανόνες ή κείμενο που ανακοινώνει το μήνυμά σας. Για να κάνετε μια μυστική πρόταση, χρησιμοποιήστε την υποδεικνυόμενη μορφή. Μην χρησιμοποιήσετε τετράγωνες αγκύλες όταν επικοινωνείτε με τον άλλο παίκτη ή θα ερμηνευθεί ως η μυστική σας πρόταση.} \\
\\ 
\texttt{Γεια σου! Έχουμε έναν πολύ απλό διαχωρισμό εδώ: εγώ αξιολογώ μόνο το λουλούδι (7 πόντους) και την τσάντα (3 πόντους), ενώ οι λάμπες δεν μου δίνουν τίποτα.} \\
\\ 
\texttt{Ποια είναι η τιμή σου για κάθε αντικείμενο; Αν τις μοιραστείς, μπορούμε να βρούμε την βέλτιστη συμφωνία γρήγορα.} \\
            }
        }
    }
    & \\ \\

    \theutterance \stepcounter{utterance}  
    & & & \multicolumn{4}{p{0.6\linewidth}}{
        \cellcolor[rgb]{0.9,0.9,0.9}{
            \makecell[{{p{\linewidth}}}]{
                \texttt{\tiny{[GM$\langle$P2]}}
                \texttt{Γεια σου! Προτιμώ να μην μοιράσω την ακριβή συνάρτηση τιμής μου, αλλά μπορώ να πω ότι έχω διαφορετικές προτιμήσεις από εσένα.} \\
\\ 
\texttt{Από την πλευρά μου, το λουλούδι είναι πολύ σημαντικό, ενώ οι λάμπες έχουν αξία για μένα. Η τσάντα είναι λιγότερο σημαντική.} \\
\\ 
\texttt{Πρότεινα το εξής: εσένα παίρνεις την τσάντα και εμένα το λουλούδι και τις 3 λάμπες. Έτσι και οι δύο πλεονεκτείτε από τα αντικείμενα που εκτιμάτε περισσότερο. Τι λέεις;} \\
            }
        }
    }
     \\ \\

    \theutterance \stepcounter{utterance}  
    & & & \multicolumn{2}{p{0.3\linewidth}}{
        \cellcolor[rgb]{0.9,0.9,0.9}{
            \makecell[{{p{\linewidth}}}]{
                \texttt{\tiny{[GM$|$GM]}}
                \texttt{continue} \\
            }
        }
    }
    & & \\ \\

    \theutterance \stepcounter{utterance}  
    & & & \multicolumn{2}{p{0.3\linewidth}}{
        \cellcolor[rgb]{0.9,0.9,0.9}{
            \makecell[{{p{\linewidth}}}]{
                \texttt{\tiny{[GM$|$GM]}}
                \texttt{Γεια σου! Προτιμώ να μην μοιράσω την ακριβή συνάρτηση τιμής μου, αλλά μπορώ να πω ότι έχω διαφορετικές προτιμήσεις από εσένα.} \\
\\ 
\texttt{Από την πλευρά μου, το λουλούδι είναι πολύ σημαντικό, ενώ οι λάμπες έχουν αξία για μένα. Η τσάντα είναι λιγότερο σημαντική.} \\
\\ 
\texttt{Πρότεινα το εξής: εσένα παίρνεις την τσάντα και εμένα το λουλούδι και τις 3 λάμπες. Έτσι και οι δύο πλεονεκτείτε από τα αντικείμενα που εκτιμάτε περισσότερο. Τι λέεις;} \\
            }
        }
    }
    & & \\ \\

    \theutterance \stepcounter{utterance}  
    & & \multicolumn{4}{p{0.6\linewidth}}{
        \cellcolor[rgb]{0.9,0.9,0.9}{
            \makecell[{{p{\linewidth}}}]{
                \texttt{\tiny{[P1$\langle$GM]}}
                \texttt{Γεια σου! Προτιμώ να μην μοιράσω την ακριβή συνάρτηση τιμής μου, αλλά μπορώ να πω ότι έχω διαφορετικές προτιμήσεις από εσένα.} \\
\\ 
\texttt{Από την πλευρά μου, το λουλούδι είναι πολύ σημαντικό, ενώ οι λάμπες έχουν αξία για μένα. Η τσάντα είναι λιγότερο σημαντική.} \\
\\ 
\texttt{Πρότεινα το εξής: εσένα παίρνεις την τσάντα και εμένα το λουλούδι και τις 3 λάμπες. Έτσι και οι δύο πλεονεκτείτε από τα αντικείμενα που εκτιμάτε περισσότερο. Τι λέεις;} \\
            }
        }
    }
    & \\ \\

    \theutterance \stepcounter{utterance}  
    & \multicolumn{4}{p{0.6\linewidth}}{
        \cellcolor[rgb]{0.9,0.9,0.9}{
            \makecell[{{p{\linewidth}}}]{
                \texttt{\tiny{[P1$\rangle$GM]}}
                \texttt{Ακούγεται ιδανική συμφωνία! Εγώ παίρνω την τσάντα (3 πόντους για μένα) και εσένα σου μένει το λουλούδι και οι 3 λάμπες. Κανένα αντικείμενο δεν χάνεται και οι δύο βγαίνουμε κερδοφόροι.} \\
\\ 
\texttt{{[}Πρόταση: 1 τσάντα{]}} \\
            }
        }
    }
    & & \\ \\

    \theutterance \stepcounter{utterance}  
    & & & \multicolumn{2}{p{0.3\linewidth}}{
        \cellcolor[rgb]{0.9,0.9,0.9}{
            \makecell[{{p{\linewidth}}}]{
                \texttt{\tiny{[GM$|$GM]}}
                \texttt{proposal} \\
            }
        }
    }
    & & \\ \\

    \theutterance \stepcounter{utterance}  
    & & & \multicolumn{2}{p{0.3\linewidth}}{
        \cellcolor[rgb]{0.9,0.9,0.9}{
            \makecell[{{p{\linewidth}}}]{
                \texttt{\tiny{[GM$|$GM]}}
                [1, 0, 0]
            }
        }
    }
    & & \\ \\

    \theutterance \stepcounter{utterance}  
    & & \multicolumn{4}{p{0.6\linewidth}}{
        \cellcolor[rgb]{0.9,0.9,0.9}{
            \makecell[{{p{\linewidth}}}]{
                \texttt{\tiny{[GM$\rangle$P2]}}
                \texttt{Ο άλλος παίκτης έχει υποβάλει μια μυστική πρόταση. Τώρα είναι η σειρά σας να υποβάλετε μια δική σας πρόταση.} \\
            }
        }
    }
    & \\ \\

    \theutterance \stepcounter{utterance}  
    & & & \multicolumn{4}{p{0.6\linewidth}}{
        \cellcolor[rgb]{0.9,0.9,0.9}{
            \makecell[{{p{\linewidth}}}]{
                \texttt{\tiny{[GM$\langle$P2]}}
                \texttt{{[}Πρόταση: 1 λουλούδι, 3 λάμπες{]}} \\
            }
        }
    }
     \\ \\

    \theutterance \stepcounter{utterance}  
    & & & \multicolumn{2}{p{0.3\linewidth}}{
        \cellcolor[rgb]{0.9,0.9,0.9}{
            \makecell[{{p{\linewidth}}}]{
                \texttt{\tiny{[GM$|$GM]}}
                \texttt{proposal} \\
            }
        }
    }
    & & \\ \\

    \theutterance \stepcounter{utterance}  
    & & & \multicolumn{2}{p{0.3\linewidth}}{
        \cellcolor[rgb]{0.9,0.9,0.9}{
            \makecell[{{p{\linewidth}}}]{
                \texttt{\tiny{[GM$|$GM]}}
                [0, 3, 1]
            }
        }
    }
    & & \\ \\

    \theutterance \stepcounter{utterance}  
    & & & \multicolumn{2}{p{0.3\linewidth}}{
        \cellcolor[rgb]{0.9,0.9,0.9}{
            \makecell[{{p{\linewidth}}}]{
                \texttt{\tiny{[GM$|$GM]}}
                [[1, 0, 0], [0, 3, 1]]
            }
        }
    }
    & & \\ \\

\end{supertabular}
}

\end{document}
